\documentclass[12pt]{article}
\usepackage[margin=0.8in]{geometry}
\usepackage{amsmath}
\usepackage{amssymb}
\usepackage{xcolor}
\usepackage{enumitem}

\title{\LARGE Variance Rules - Simple Guide}
\author{ECO 310}
\date{}

\setlength{\parskip}{0.5em}

\begin{document}

\maketitle

\section*{\Large The 3 Most Important Rules}

\subsection*{\large Rule 1: Constants Get Squared}

\begin{center}
\colorbox{yellow!30}{
\Large
$\text{var}(aX) = a^2 \cdot \text{var}(X)$
}
\end{center}

\textbf{In words:} When you multiply a random variable by a number, square that number.

\vspace{1em}

\textbf{Example:}
\begin{itemize}[leftmargin=2cm]
    \item If $\text{var}(X) = 4$
    \item Then $\text{var}(3X) = 3^2 \times 4 = 9 \times 4 = 36$
\end{itemize}

\vspace{1em}

\textbf{Common mistake:} $\text{var}(3X) \neq 3 \times 4$ ~~~ You must square the 3!

\vspace{2em}

\subsection*{\large Rule 2: Adding Constants Does Nothing}

\begin{center}
\colorbox{yellow!30}{
\Large
$\text{var}(X + c) = \text{var}(X)$
}
\end{center}

\textbf{In words:} Adding or subtracting a constant doesn't change variance.

\vspace{1em}

\textbf{Example:}
\begin{itemize}[leftmargin=2cm]
    \item If $\text{var}(X) = 9$
    \item Then $\text{var}(X + 100) = 9$
    \item And $\text{var}(X - 50) = 9$
\end{itemize}

\vspace{1em}

\textbf{Why?} Variance measures spread. Shifting everything by the same amount doesn't change the spread.

\newpage

\subsection*{\large Rule 3: The Big Formula (Most Important!)}

\begin{center}
\colorbox{red!20}{
\LARGE
$\text{var}(aX + bY) = a^2\text{var}(X) + b^2\text{var}(Y) + 2ab\text{cov}(X,Y)$
}
\end{center}

\vspace{1em}

\textbf{Special case - If X and Y are independent:}

\begin{center}
\colorbox{green!20}{
\LARGE
$\text{var}(aX + bY) = a^2\text{var}(X) + b^2\text{var}(Y)$
}
\end{center}

(The covariance term disappears!)

\vspace{2em}

\section*{\Large Step-by-Step Example}

\subsection*{Problem:}

You invest in two assets:
\begin{itemize}[leftmargin=2cm]
    \item Asset A: $\text{var}(A) = 0.25$
    \item Asset B: risk-free, so $\text{var}(B) = 0$
    \item You put fraction $\alpha = 0.5$ in A, and $(1-\alpha) = 0.5$ in B
\end{itemize}

\textbf{Find:} $\text{var}(\alpha A + (1-\alpha)B)$

\vspace{1em}

\subsection*{Solution:}

\textbf{Step 1:} Write down the formula
\[
\text{var}(\alpha A + (1-\alpha)B) = \alpha^2\text{var}(A) + (1-\alpha)^2\text{var}(B) + 2\alpha(1-\alpha)\text{cov}(A,B)
\]

\textbf{Step 2:} Since B is risk-free, $\text{var}(B) = 0$ and $\text{cov}(A,B) = 0$
\[
\text{var}(\alpha A + (1-\alpha)B) = \alpha^2\text{var}(A)
\]

\textbf{Step 3:} Plug in numbers
\[
= (0.5)^2 \times 0.25 = 0.25 \times 0.25 = 0.0625
\]

\textbf{Answer:} Portfolio variance = 0.0625

\newpage

\section*{\Large Common Patterns}

\subsection*{\large Pattern 1: Just Adding Two Variables}

\begin{center}
\Large
$\text{var}(X + Y)$ = ?
\end{center}

\vspace{1em}

\textbf{If independent:}
\[
\text{var}(X + Y) = \text{var}(X) + \text{var}(Y)
\]

\textbf{If NOT independent:}
\[
\text{var}(X + Y) = \text{var}(X) + \text{var}(Y) + 2\text{cov}(X,Y)
\]

\vspace{2em}

\subsection*{\large Pattern 2: Subtracting Variables}

\begin{center}
\Large
$\text{var}(X - Y)$ = ?
\end{center}

\vspace{1em}

\textbf{If independent:}
\[
\text{var}(X - Y) = \text{var}(X) + \text{var}(Y)
\]

\textbf{Notice:} It's PLUS var(Y), not minus!

\textbf{If NOT independent:}
\[
\text{var}(X - Y) = \text{var}(X) + \text{var}(Y) - 2\text{cov}(X,Y)
\]

\newpage

\section*{\Large Quick Reference}

\begin{center}
\begin{tabular}{|p{5cm}|p{8cm}|}
\hline
\textbf{\large What You Have} & \textbf{\large Variance} \\
\hline
\hline
A constant $c$ & $0$ \\
\hline
$X + 5$ & $\text{var}(X)$ \\
\hline
$3X$ & $9 \cdot \text{var}(X)$ \\
\hline
$X + Y$ (independent) & $\text{var}(X) + \text{var}(Y)$ \\
\hline
$X - Y$ (independent) & $\text{var}(X) + \text{var}(Y)$ \\
\hline
$2X + 3Y$ (independent) & $4\text{var}(X) + 9\text{var}(Y)$ \\
\hline
\end{tabular}
\end{center}

\vspace{2em}

\section*{\Large The 3 Biggest Mistakes}

\subsection*{Mistake 1: Forgetting to Square}

\textbf{WRONG:} $\text{var}(5X) = 5 \cdot \text{var}(X)$

\textbf{RIGHT:} $\text{var}(5X) = 25 \cdot \text{var}(X)$

\vspace{1em}

\subsection*{Mistake 2: Subtracting Variances}

\textbf{WRONG:} $\text{var}(X - Y) = \text{var}(X) - \text{var}(Y)$

\textbf{RIGHT:} $\text{var}(X - Y) = \text{var}(X) + \text{var}(Y)$ (if independent)

\vspace{1em}

\subsection*{Mistake 3: Assuming Independence}

\textbf{WRONG:} Always using $\text{var}(X + Y) = \text{var}(X) + \text{var}(Y)$

\textbf{RIGHT:} Only works if X and Y are independent! Otherwise need the covariance term.

\newpage

\section*{\Large Practice Problems}

Given: $\text{var}(X) = 16$, $\text{var}(Y) = 9$, and X and Y are independent.

\vspace{1em}

\textbf{Problem 1:} Find $\text{var}(2X)$

\textbf{Solution:} $\text{var}(2X) = 2^2 \times 16 = 4 \times 16 = 64$

\vspace{2em}

\textbf{Problem 2:} Find $\text{var}(X + 7)$

\textbf{Solution:} $\text{var}(X + 7) = \text{var}(X) = 16$

\vspace{2em}

\textbf{Problem 3:} Find $\text{var}(X + Y)$

\textbf{Solution:} Since independent: $\text{var}(X + Y) = 16 + 9 = 25$

\vspace{2em}

\textbf{Problem 4:} Find $\text{var}(3X - 2Y)$

\textbf{Solution:}
\begin{align*}
\text{var}(3X - 2Y) &= 3^2\text{var}(X) + (-2)^2\text{var}(Y) \\
&= 9 \times 16 + 4 \times 9 \\
&= 144 + 36 \\
&= 180
\end{align*}

\vspace{2em}

\textbf{Problem 5:} Find $\text{var}(0.5X + 0.5Y)$

\textbf{Solution:}
\begin{align*}
\text{var}(0.5X + 0.5Y) &= (0.5)^2\text{var}(X) + (0.5)^2\text{var}(Y) \\
&= 0.25 \times 16 + 0.25 \times 9 \\
&= 4 + 2.25 \\
&= 6.25
\end{align*}

\newpage

\section*{\Large For the Midterm}

\subsection*{What You Need to Remember:}

\begin{enumerate}[leftmargin=2cm, itemsep=1em]
    \item \textbf{Square the constants:} $\text{var}(aX) = a^2\text{var}(X)$

    \item \textbf{The big formula:}
    \[
    \text{var}(aX + bY) = a^2\text{var}(X) + b^2\text{var}(Y) + 2ab\text{cov}(X,Y)
    \]

    \item \textbf{If independent:} Drop the covariance term

    \item \textbf{Risk-free asset:} Has variance = 0

    \item \textbf{Standard deviation:} $\sigma = \sqrt{\text{variance}}$
\end{enumerate}

\vspace{2em}

\subsection*{The Portfolio Problem Pattern:}

If you invest $\alpha$ in risky asset A and $(1-\alpha)$ in risk-free asset B:

\begin{center}
\colorbox{blue!10}{
\large
$\text{var}(\text{portfolio}) = \alpha^2 \text{var}(A)$
}
\end{center}

\vspace{1em}

And standard deviation:

\begin{center}
\colorbox{blue!10}{
\large
$\sigma_{\text{portfolio}} = |\alpha| \cdot \sigma_A$
}
\end{center}

\newpage

\section*{\LARGE MC Problem 4: Full Solution}

\subsection*{\large The Problem}

There are two assets available:
\begin{itemize}[leftmargin=2cm]
    \item \textbf{Risk-free asset:} 10\% return rate
    \item \textbf{Risky asset:} 20\% expected return rate and variance = $\frac{1}{4}$
\end{itemize}

\vspace{1em}

You divide your money between the two assets so that the \textbf{standard deviation} of your overall investment is $\frac{1}{4}$.

\vspace{1em}

\textbf{Question:} What is your overall expected return rate?

\begin{enumerate}[label=(\alph*), leftmargin=2cm]
    \item 0.15
    \item 0.1
    \item 0.2
    \item None of these is correct
\end{enumerate}

\vspace{2em}

\subsection*{\large Step 1: Set Up the Portfolio}

Let $\alpha$ = fraction of money in the risky asset

Then $(1-\alpha)$ = fraction in the risk-free asset

\vspace{1em}

Your portfolio return is:
\[
R = \alpha R_{\text{risky}} + (1-\alpha) R_{\text{risk-free}}
\]

\vspace{2em}

\subsection*{\large Step 2: Find the Variance of the Risky Asset's Standard Deviation}

We're given: $\text{var}(R_{\text{risky}}) = \frac{1}{4}$

Therefore:
\[
\sigma_{\text{risky}} = \sqrt{\frac{1}{4}} = \frac{1}{2}
\]

\newpage

\subsection*{\large Step 3: Find Portfolio Variance}

Since the risk-free asset has NO variance:
\[
\text{var}(R_{\text{risk-free}}) = 0
\]

Using our formula:
\begin{align*}
\text{var}(\text{portfolio}) &= \text{var}(\alpha R_{\text{risky}} + (1-\alpha) R_{\text{risk-free}}) \\
&= \alpha^2 \text{var}(R_{\text{risky}}) + (1-\alpha)^2 \text{var}(R_{\text{risk-free}}) \\
&= \alpha^2 \times \frac{1}{4} + (1-\alpha)^2 \times 0 \\
&= \alpha^2 \times \frac{1}{4}
\end{align*}

\vspace{1em}

Therefore, portfolio standard deviation is:
\[
\sigma_{\text{portfolio}} = \sqrt{\alpha^2 \times \frac{1}{4}} = \alpha \times \frac{1}{2}
\]

\vspace{2em}

\subsection*{\large Step 4: Use the Given Information}

We're told the portfolio standard deviation is $\frac{1}{4}$:

\[
\alpha \times \frac{1}{2} = \frac{1}{4}
\]

Solve for $\alpha$:
\[
\alpha = \frac{1/4}{1/2} = \frac{1}{4} \times \frac{2}{1} = \frac{1}{2}
\]

\vspace{1em}

\begin{center}
\colorbox{yellow!30}{
\large
So you invest \textbf{50\%} in the risky asset and \textbf{50\%} in the risk-free asset.
}
\end{center}

\newpage

\subsection*{\large Step 5: Calculate Expected Return}

Expected return is a \textbf{weighted average}:

\begin{align*}
E[R] &= \alpha \times E[R_{\text{risky}}] + (1-\alpha) \times E[R_{\text{risk-free}}] \\
&= \frac{1}{2} \times 0.20 + \frac{1}{2} \times 0.10 \\
&= 0.10 + 0.05 \\
&= 0.15
\end{align*}

\vspace{2em}

\begin{center}
\colorbox{green!20}{
\LARGE
\textbf{Answer: (a) 0.15 or 15\%}
}
\end{center}

\vspace{2em}

\subsection*{\large Key Takeaways from MC4}

\begin{enumerate}[leftmargin=2cm, itemsep=1em]
    \item \textbf{Risk-free means zero variance:} $\text{var}(R_{\text{risk-free}}) = 0$

    \item \textbf{Portfolio with one risk-free asset:}
    \[
    \sigma_{\text{portfolio}} = \alpha \cdot \sigma_{\text{risky}}
    \]
    (Just multiply the fraction by the risky asset's std dev)

    \item \textbf{Standard deviation is the square root of variance:}
    \[
    \sigma = \sqrt{\text{var}}
    \]

    \item \textbf{Expected return is a weighted average:}
    \[
    E[R] = \alpha \times r_1 + (1-\alpha) \times r_2
    \]
    (No need for variance here, just multiply fractions by returns)
\end{enumerate}

\vspace{2em}

\subsection*{\large Why This Works}

When you mix a risky asset with a risk-free asset:
\begin{itemize}[leftmargin=2cm]
    \item The risk-free part contributes \textbf{nothing to variance} (it's constant)
    \item Only the risky part adds variance
    \item So portfolio variance = $\alpha^2 \times$ (risky variance)
    \item The $\alpha$ controls how much risk you take on
\end{itemize}

\end{document}
